%%% =========================================================================
%%% The Explainability Gap: A Mathematical Limit on AI Regulation
%%% Nature Comment / Correspondence (2-3 pages)
%%% =========================================================================
\documentclass[11pt]{article}
\usepackage[margin=1in]{geometry}
\usepackage[utf8]{inputenc}
\usepackage[T1]{fontenc}
\usepackage{amsmath,amssymb}
\usepackage{booktabs}
\usepackage{hyperref}
\usepackage{natbib}

\title{The Explainability Gap:\\
  Why AI Regulation Demands the Mathematically Impossible---and What To Do Instead}

\author{Drake Caraker, Bryan Arnold, David Rhoads}

\date{April 2026}

\begin{document}

\maketitle

% =========================================================================
% Note: This is a companion piece to "The Limits of Explanation" (main paper).
% Target: Nature Comment, Nature Machine Intelligence Comment, or Science Policy Forum.
% Length: 1500-2000 words.
% =========================================================================

\section*{The regulatory demand}

The EU AI Act (Regulation 2024/1689) requires that high-risk AI systems be
``designed and developed in such a way as to ensure that their operation is
sufficiently transparent to enable deployers to interpret a system's output
and use it appropriately'' (Article~13).  Article~86 establishes a ``right
to explanation'' for individuals affected by AI decisions.  The US Federal
Reserve's SR~11-7 requires ``developmental evidence'' documenting how model
outputs arise.  The FDA's guidance on AI/ML-based software demands
``transparency'' in clinical decision support.

These regulations share an implicit assumption: that a sufficiently careful
explanation of an AI system can be simultaneously \emph{correct} (faithful to
the model's actual reasoning), \emph{reproducible} (stable across equivalent
model instances), and \emph{complete} (decisive about every relevant
distinction).

\section*{The mathematical limit}

We have proved~\citep{caraker2026universal} that this assumption is false.
For any AI system where multiple model instances achieve the same predictive
performance---the \emph{Rashomon property}, which holds for virtually all
production ML systems with more parameters than
constraints~\citep{breiman2001statistical,rudin2024amazing}---no explanation
can be simultaneously faithful, stable, and decisive.  The proof is four
steps, verified in the Lean~4 proof assistant (104~files, 463~theorems,
0~unproved goals), and applies uniformly to all major explanation methods:
SHAP, attention maps, counterfactual explanations, concept probes, saliency
maps, causal discovery, and mechanistic interpretability.

The impossibility is not a deficiency of current methods.  It is a
structural consequence of underspecification---a mathematical ceiling that no
future method can exceed.

\section*{What regulations actually require}

Mapping the regulatory language onto our formal framework:

\begin{center}
\begin{tabular}{@{}p{5cm}p{3cm}p{3cm}@{}}
\toprule
\textbf{Regulatory requirement} & \textbf{Formal property} & \textbf{Achievable alone?} \\
\midrule
``Interpret a system's output'' & Faithfulness & Yes \\
``Use it appropriately'' across deployments & Stability & Yes \\
``Right to explanation'' (complete) & Decisiveness & Yes \\
All three simultaneously & F + S + D & \textbf{No} \\
\bottomrule
\end{tabular}
\end{center}

Each property is achievable individually or in pairs.  The impossibility
arises only when all three are demanded simultaneously for underspecified
systems.  The bilemma strengthens this: for systems with binary explanation
spaces (e.g., ``feature is important'' vs.\ ``not important''), even
faithful + stable alone is impossible without enriching the explanation space
with an indeterminate option.

\section*{The practical impact}

How severe is this?  We audited five clinical and financial datasets commonly
used in published SHAP analyses (Breast Cancer Wisconsin, Heart Disease,
Pima Diabetes, German Credit, Adult Income) using TreeSHAP with 30
independently trained XGBoost models per dataset.

\textbf{28\% of top-5 SHAP feature ranking comparisons are structurally
unreliable}---meaning that retraining the model with a different random seed
produces the opposite ranking with $>$30\% probability.  On the Breast
Cancer dataset, the comparison between the two most important features
(worst perimeter vs.\ worst concave points) is a coin flip (SNR = 0.17,
43\% flip rate).  On Adult Income, 60\% of top-5 comparisons are unreliable.

These are not errors by the analysts.  They are structural consequences of
the Rashomon property.  No existing regulatory framework diagnoses or
accounts for this instability.

\section*{The constructive solution}

The impossibility theorem does not counsel despair.  It provides a
constructive resolution: \emph{$G$-invariant projection}, which aggregates
explanations across the Rashomon set.  The concrete implementation---the
SAGE algorithm (Symmetry-Aware Group Estimation)---identifies which
explanation queries are structurally reliable (signal-to-noise ratio $> 2$)
and which are structurally unreliable (SNR $< 0.5$), reducing false
discovery from 57\% to $<$5\%.

The coverage conflict bridge validates this quantitatively: across 15
datasets, the theoretical predictor (coverage conflict degree, derived from
the impossibility theorem) correlates with observed instability at
$\rho = 0.961$ ($p < 0.0001$).

\section*{Recommendations for regulators}

We propose that AI explainability regulations adopt the following standard:

\begin{enumerate}
  \item \textbf{Require stability-annotated explanations.}  Every feature
    ranking or explanation should report which comparisons are structurally
    reliable (SNR $> 2$), marginal ($0.5 < \text{SNR} < 2$), or unreliable
    (SNR $< 0.5$).

  \item \textbf{Mandate the Rashomon diagnostic.}  Before requiring
    explanations, assess whether the system has the Rashomon property
    (multiple equivalent models exist).  If it does, require
    Rashomon-aware explanation methods.

  \item \textbf{Replace ``complete explanation'' with ``best achievable
    explanation.''}  The optimal tradeoff is known:
    $G$-invariant aggregation maximizes faithfulness among stable methods.
    Require this optimality, not the impossible triple.

  \item \textbf{Require explicit disclosure of structural indeterminacy.}
    When explanation queries are unreliable, the system should report ``this
    comparison is structurally indeterminate'' rather than presenting an
    arbitrary answer as definitive.
\end{enumerate}

These recommendations are constructive, not obstructive.  They replace an
impossible demand (complete, stable, faithful explanations) with an achievable
standard (stability-annotated, Rashomon-aware, optimally aggregated
explanations) that is provably the best any method can do.

\section*{The precedent}

This parallels Arrow's impossibility theorem~\citep{arrow1951social} in
social choice theory.  Arrow proved that no voting system can simultaneously
satisfy unanimity, independence of irrelevant alternatives, and
non-dictatorship.  The response was not to abandon voting but to
characterize the achievable tradeoffs and design mechanisms accordingly.
Our result serves the same function for AI explainability: it replaces an
implicit assumption of possibility with a precise characterization of what
is and is not achievable, enabling regulation grounded in mathematical
reality rather than unrealizable aspiration.

\bigskip
\noindent\textit{The companion paper ``The Limits of Explanation'' provides
the full formal framework, Lean~4 verification, and empirical validation
across 23 domains.}

\bibliographystyle{plain}
\bibliography{references}

\end{document}
